\documentclass[11pt]{article}
% Shared preamble for all daily exercises
% Update this file to change settings (padding, parindent, etc.) for all exercises

\usepackage[margin=0.75in]{geometry}
\usepackage{amsmath}
\usepackage{amssymb}

% Custom settings
\setlength{\parindent}{0pt}
\setlength{\parskip}{0.2cm}

\title{Daily Exercise}
\author{Dhruman Gupta}
\date{28th March}

\begin{document}

% Common header for daily exercises
% This file can be included in other LaTeX documents

\maketitle

\noindent
\textbf{Course:} CS-3330, Theory of Computation

\vspace{0.2cm}

\noindent
\textbf{Question:}\noindent 
Consider the following decision problem: given a Python program $P$ and a line number $\ell$, does there exist an input on which $P$ reaches line $\ell$ during execution? Show that this problem is undecidable.

\textbf{Answer:}

We will prove this by reduction from the halting problem.

Let $M$ be a program and let $w$ be an input. We will construct a Python program $Q$ and a line number $\ell$ such that $Q$ reaches line $\ell$ on some input if and only if $M$ halts on $w$.

The program $Q$ works as follows:

\begin{itemize}
    \item It takes an arbitrary input, but ignores it.
    \item It simulates $M$ on input $w$.
    \item If the simulation of $M$ on $w$ halts, then $Q$ proceeds to line $\ell$.
    \item If the simulation does not halt, then $Q$ loops forever before reaching line $\ell$.
\end{itemize}

Now, if $M$ halts on $w$, then the simulation inside $Q$ eventually finishes, and so $Q$ reaches line $\ell$. Hence there exists an input on which $Q$ reaches line $\ell$.

On the other hand, if $M$ does not halt on $w$, then the simulation inside $Q$ runs forever, so line $\ell$ is never reached on any input.

Therefore,
\[
Q \text{ reaches line } \ell \text{ on some input } \iff M \text{ halts on } w.
\]

So if we had a decider for the given problem, then we could decide the halting problem by applying it to $Q$ and $\ell$. But the halting problem is undecidable. This is a contradiction.

Hence the problem of deciding whether there exists an input on which a Python program reaches a given line is undecidable.

\end{document}
